\chapter{Testing: ERT, Buttercup, and why the tests matter}
\label{ch:testing}
\minitoc\label{testing-ert-buttercup-and-why-the-tests-matter}
\index{ERT}\index{Buttercup}\index{testing}\index{regression testing}

% ── Opening signpost ──────────────────────────────────────────────
Every performance change in Part~\ref{part:startup} carries a risk:
the optimisation works, but a feature silently breaks.
An Emacs config is a program, and programs need tests.  This chapter
covers ERT and Buttercup (the two testing frameworks available in
Emacs) and explains the specific regression vectors that the
configuration's test suite guards against.  Understanding these
patterns lets you add tests for any new feature you introduce.

\section{Why test an Emacs config?}\label{why-test-an-emacs-config}

An Emacs config is not a static document; it is a program that runs on startup. Performance regressions are introduced silently: a package update changes a default, a new \texttt{use-package} declaration accidentally uses \texttt{:config} instead of \texttt{:init}, or a global mode resets a variable. Without tests, these regressions are discovered only when startup becomes painfully slow weeks later, with no git bisect breadcrumb.

\section{ERT (Emacs Lisp Regression Testing)}\label{ert-emacs-lisp-regression-testing}

\texttt{(info "(ert) Top")}

ERT is built into Emacs. Tests are defined with \texttt{ert-deftest} and run with \texttt{M-x ert} or from the command line.

\subsection{Anatomy of a test}\label{anatomy-of-a-test}

\begin{minted}{elisp}
(ert-deftest test/gc-threshold-matches-early-init ()
  "gc-cons-threshold must equal my/gc-cons-threshold after startup."
  (should (boundp 'my/gc-cons-threshold))
  (should (= gc-cons-threshold my/gc-cons-threshold)))
\end{minted}

\begin{itemize}
\item \texttt{should} is the assertion macro. Test fails if the body evaluates to \texttt{nil}.
\item \texttt{should-not} asserts the body is \texttt{nil}.
\item \texttt{should-error} asserts the body signals an error.
\end{itemize}

\subsection{Running tests from the command line}\label{running-tests-from-the-command-line}

\begin{minted}{bash}
gtimeout 30 emacs --batch -l ~/.emacs.d/init.el \
  -l tests/test-performance.el \
  -f ert-run-tests-batch-and-exit
\end{minted}

\texttt{-\/-batch} means no GUI, output to stdout/stderr, exit on completion. The exit code is 0 if all tests pass, non-zero otherwise.

\subsection{The Emacs-under-test problem}\label{the-emacs-under-test-problem}

When you run \texttt{emacs -\/-batch -l init.el -l tests.el}, the tests run in a fully initialised Emacs with all your packages loaded. This is both an advantage (you're testing the real config) and a limitation:

\begin{itemize}
\item Tests that check mode activation (e.g., \texttt{savehist-mode}) need the corresponding hooks to have fired. In \texttt{-\/-batch} mode, \texttt{emacs-startup-hook} fires, but \texttt{window-setup-hook} does \textbf{not} (there's no window to set up).
\item Tests that use \texttt{sit-for} (e.g., waiting for an idle timer) won't work in \texttt{-\/-batch} because there's no event loop to trigger idle.
\end{itemize}

For idle-timer tests, you need to run them interactively or use a workaround:

\begin{minted}{elisp}
;; Force idle timers to fire in batch mode:
(let ((timers (copy-sequence timer-idle-list)))
  (dolist (timer timers)
    (timer-event-handler timer)))
\end{minted}

\section{Buttercup (BDD-style testing)}\label{buttercup-bdd-style-testing}

Buttercup provides \texttt{describe} / \texttt{it} / \texttt{expect} syntax, which reads more naturally for complex test suites:

\begin{minted}{elisp}
(describe "GC configuration"
  (it "maintains the threshold set in early-init.el"
    (expect gc-cons-threshold :to-equal my/gc-cons-threshold))

  (it "has completed fewer than 20 GCs during startup"
    (expect gcs-done :to-be-less-than 20)))
\end{minted}

\subsection{Running Buttercup}\label{running-buttercup}

\begin{minted}{bash}
gtimeout 30 emacs --batch -l ~/.emacs.d/init.el \
  -l buttercup -f buttercup-run
\end{minted}

Buttercup discovers test files by convention (files matching \texttt{*-test.el} or \texttt{test-*.el}) or you can specify them explicitly.

\section{What the suggested tests actually verify}\label{what-the-suggested-tests-actually-verify}

The tests in the fix log are not arbitrary: each one guards a specific regression vector:

\begin{longtable}[]{@{}ll@{}}
\toprule
Test category & What it catches \\
\midrule
GC threshold check & Another package resetting \texttt{gc-cons-threshold} \\
\texttt{gcmh-mode} not active & Accidental re-introduction of GCMH \\
Global modes active & A mode silently disabled by \texttt{always-defer} \\
Keybinding checks & A package overriding your binding \\
\texttt{general.el} not loaded & Accidental re-addition of the package \\
\texttt{editorconfig-mode} active & Regression of the "dead \texttt{:config}" bug \\
Dashboard buffer exists & Dashboard broken again by defer changes \\
Server running after idle & \texttt{server-start} removed or broken \\
\bottomrule
\end{longtable}

\subsection{The "feature still works" pattern}\label{the-feature-still-works-pattern}

For every performance change that touches user-facing functionality, include a test that exercises the feature:

\begin{minted}{elisp}
(ert-deftest test/git-gutter-keybindings ()
  "git-gutter keybindings should remain accessible."
  (should (eq (keymap-lookup global-map "C-c g n")
              'git-gutter:next-hunk)))
\end{minted}

This catches the case where you defer git-gutter but accidentally break its keybindings in the process.

\section{Limitations}\label{limitations}

\begin{itemize}
\item ERT/Buttercup test the config \textbf{as loaded}. They cannot test the load \textbf{process} itself (timing, order). For that, use the profiling techniques in Chapter~\ref{ch:profiling}.
\item \texttt{-\/-batch} mode does not fire \texttt{window-setup-hook}. Tests for dashboard, doom-modeline, or other window-dependent features must be run interactively or must force the hooks.
\item Tests that rely on external state (Git repos, LSP servers, network) are inherently flaky. Guard them with \texttt{:condition} or \texttt{skip-unless}:
\end{itemize}

\begin{minted}{elisp}
(ert-deftest test/git-gutter-activates () (skip-unless
  (bound-and-true-p global-git-gutter-mode)) (should ...))
\end{minted}

% ── Mental model summary ──────────────────────────────────────────
\section{Mental model}
Emacs ships ERT (\texttt{ert-deftest}) as its built-in test framework,
while Buttercup (\texttt{describe}/\texttt{it}/\texttt{expect}) layers
BDD-style syntax on top.  Both frameworks run inside a fully initialised
\texttt{--batch} Emacs, which means the tests verify the real
configuration rather than a mock, but \texttt{window-setup-hook} does
not fire in batch mode, so any feature that depends on it needs a
workaround.  Each test should guard a specific regression vector: a GC
threshold that gets silently reset, a \texttt{:config} block that never
executes, or a keybinding that another package overrides.  The
complementary ``feature still works'' pattern exercises the user-facing
command after every performance change, ensuring that optimisation does
not break functionality.  Where a test depends on external state (a Git
repository, an LSP server), \texttt{skip-unless} gates execution so the
absence of that state produces a skip rather than a false failure.
